\documentclass[11pt,a4paper]{report}
\usepackage{marvosym}

\assignment{1}
\group{...}
\students{..........}{..........}

\begin{document}

\maketitle

Answer to the questions by adding text in the boxes. 
You may answer in either \textbf{French or English}. 
Do not modify anything else in the template.  
The size of the boxes indicate the place you \textbf{can} use, but you \textbf{need not} use all of it (it is not an indication of any expected length of answer). 
\textbf{Be as concise as possible! A good short answer is better than a lot of nonsense!}
%\bigskip

\section{Understanding a CP model (2 points)}

\begin{enumerate}
    \item Answer correctly to the questions 1-4. \textbf{(0.5pt)}
    You can help yourself by looking
    at \url{https://www.pycsp.org/documentation/constraints/} to find the definition of the different constraints.
    \item Answer correctly to the question 5. \textbf{(0.5pt)}
    \item Add the needed constraint to the model and submit it on INGInious. \textbf{(1pt)}
\end{enumerate}

\newpage
\section{Minesweeper Model (8 points)}

\begin{enumerate}
    \item Implement your CP model in the file \texttt{minesweeper.py} and submit it on INGInious. (\textbf{4 pts})
    \item Describe your CP model: First describe the variables of your CP model. 
    For each set of variables, explain what they represent, give their domains and how many variables are in this set. 
    Then, list the constraints of your CP model. 
    For each constraint explain what it enforces, and on which variables it is applied. (\textbf{4 pts})
\end{enumerate}

\begin{answers}[10cm]
\textbf{Variables :}
\begin{itemize}
    \item \textbf{Représentation :} Un tableau 2D de variables, par exemple \texttt{mines[i][j]},  représentant la grille du jeu.
    \item \textbf{Signification :} Chaque variable représente la présence (1) ou l'absence (0) d'une mine à la ligne $i$ et colonne $j$.
    \item \textbf{Domaine :} $\{0, 1\}$ pour chaque variable.
    \item \textbf{Quantité :} $n \times m$ variables (où $n$  est le nombre de lignes et $m$ le nombre de colonnes de la grille).
\end{itemize}
\textbf{Contraintes :}
\begin{itemize}
    \item \textbf{Contraintes 1 (Indices sécurisés) :} Pour chaque case $(i,j)$ contenant un indice  ($clues[i][j] \neq -1$), on impose \texttt{mines[i][j] == 0}. Cela garantit qu'aucune mine n'est placée sur une case indice. Appliquée directementsur la variable de la case.
    \item \textbf{Contrainte 2 (Comptage des mines) :} Pour chaque case indice $(i,j)$, on impose \texttt{Sum(neighbors) == clues[i][j]}. Cela garantit que le nombreexact de mines dans les cases adjacentes correspond à l'indice. Appliquée sur l'ensemble des variables voisines valides (jusqu'à 8).
\end{itemize}
\end{answers}




\newpage
\section{Atom Placement Problem()}
\begin{enumerate}
\item Formulate the atom placement problem as a Local Search problem (problem, cost function, feasible solutions, optimal solutions). \textbf{(1 pt)}
\end{enumerate}

\begin{answers}[6cm]
% Your answer here
\end{answers}


\begin{enumerate}
\setcounter{enumi}{1}
    \item You are given a template on Moodle: \textit{atomplacement.py}. Implement your own extension of the \textit{Problem} class from \texttt{aima-python3}. Implement the \texttt{maxvalue} and \texttt{randomized maxvalue} strategies. To do so,
	you can get inspiration from the \textit{randomwalk} function in \textit{search.py}. Your program will be evaluated on 15 instances (during 1 minute) of which 5 are hidden. We expect you to solve at least 12 out the 15.
    \begin{enumerate}
        \item \texttt{maxvalue} chooses the best node (i.e., the node with maximal
        value) in the neighborhood, even if it degrades the quality of the current solution. 
        The \texttt{maxvalue} strategy should be defined in a function called
        \textit{maxvalue} with the following signature: \\
        \texttt{maxvalue(problem,limit=100)}. \textbf{(2.5 pts)}
    \end{enumerate}
\end{enumerate}

\begin{answers}[2.5 cm]
% max_value comments
pour éviter de rester bloqué dans de mauvais optima locaux, la descente a été englobée dans une boucle de "Random Restart" (15 redémarrages). À chaque essai, l'état initial est mélangé aléatoirement, et on conserve le meilleur résultat global.
\end{answers}



\begin{enumerate}
\setcounter{enumi}{1}
\begin{enumerate}
\setcounter{enumii}{1}
    \item \texttt{randomized maxvalue} chooses the next node randomly among
        the 5~best neighbors (again, even if it degrades the quality of the current solution).
        The \texttt{randomized maxvalue} strategy should be defined in a function called
        \textit{randomized\_maxvalue} with the following signature: \\
        \texttt{randomized\_maxvalue(problem,limit=100)}. \textbf{(2.5 pts)}
\end{enumerate}
\end{enumerate}

\begin{answers}[2.5 cm]
Similaire à max\_value, nous utilisons 15 redémarrages aléatoires. Lors de la recherche locale, les voisins sont triés par énergie croissante, et le prochain état est choisi par hasard parmi les 5 meilleurs pour introduire une légère diversification.
\end{answers}



\begin{enumerate}
\setcounter{enumi}{2}
\item Compare the 2 strategies implemented in the previous question and \texttt{randomwalk} defined in \textit{search.py} on the given vertex cover instances. 
    Run the strategies during 100 steps, and report, in a table, the computation
    time, the value of the best solution and the number of steps needed to reach the best result. For the
    randomized max value and the random walk, each instance should be tested 10~times to reduce
    the effects of the randomness on the result. When multiple runs of the same instance
    are executed, report the mean of the quantities. \textbf{(3 pts)}
\end{enumerate}

\begin{answers}[7cm]
% Your answer here
\begin{center}
\begin{tabular}{||l||l|l|l||l|l|l||l|l|l||}
\hline
\multirow{3}{*}{Inst.} & \multicolumn{3}{c||}{Max value} & \multicolumn{3}{c||}{Random max value} & \multicolumn{3}{c||}{Random walk} \\
\cline{2-10}
\cline{2-10}
 & T(s) & Val & NS & T(s) & Val & NS & T(s) & Val & NS\\
\hline
i01 & & & & & & & & &\\
\hline
i02 & & & & & & & & &\\
\hline
i03 & & & & & & & & &\\
\hline
i04 & & & & & & & & &\\
\hline
i05 & & & & & & & & &\\
\hline
i06 & & & & & & & & &\\
\hline
i07 & & & & & & & & &\\
\hline
i08 & & & & & & & & &\\
\hline
i09 & & & & & & & & &\\
\hline
i10 & & & & & & & & &\\
\hline
\end{tabular}
\end{center}
\textbf{NS}: Number of steps — \textbf{T}: Time — \textbf{Val}: Value of the best solution
\end{answers}



\begin{enumerate}
\setcounter{enumi}{3}
    \item \textbf{(4 pts)} Answer the following questions:
    \begin{enumerate}
        \item In one sentence, what is the best strategy? \textbf{(0.5 pt)}
    \end{enumerate}
\end{enumerate}

\begin{answers}[3.5cm]
La meilleure stratégir est \texttt{randomized\_maxvalue} (combinée à des redémarrages aléatoires), car elle  équilibre efficacement l'exploitation localeet la capacité à s'échapper des petits optima locaux.
\end{answers}


\newpage
\begin{enumerate}
\setcounter{enumi}{3}
\begin{enumerate}
\setcounter{enumii}{1}
    \item Why do you think the best strategy beats the other ones? What conclusions can be drawn on the atom placement problem ? \textbf{(1.5 pt)}
\end{enumerate}
\end{enumerate}

\begin{answers}[7.5cm]
Le problème de placement d'atomes possède un espace de recherche non convexe rempli de nombreux optima locaux. Une approche \texttt{max\_value} pure est trop gourmande (gloutonne) et se retrouve piégée dans la première  "vallée" rencontrée. À l'inverse, \texttt{random\_walk} est trop chaotique. \texttt{randomized\_maxvalue} introduit juste assez de bruit (choix parmi les 5 meilleurs) pour franchir les petites  barrières d'énergie, permettant  d'atteindre des vallées plus profondes (meilleures solutions).
\end{answers}



\begin{enumerate}
\setcounter{enumi}{3}
\begin{enumerate}
\setcounter{enumii}{2}
    \item What are the limitations of each strategy in terms of diversification 
    and intensification? \textbf{(1 pt)}
\end{enumerate}
\end{enumerate}

\begin{answers}[5cm]
\beg
\end{answers}



\begin{enumerate}
\setcounter{enumi}{3}
\begin{enumerate}
\setcounter{enumii}{2}
    \item What is the behavior of the different techniques when they fall 
    in a local optimum? \textbf{(1 pt)}
\end{enumerate}
\end{enumerate}

\begin{answers}[5cm]
\begin{itemize}
    \item \textbf{max\_value :} L'algorithme s'arrête définitivement. N'acceptant aucune dégradation de la solution, il reste coincé au fond du premier optimum local rencontré.
    \item \textbf{randomized\_maxvalue :} Peut potentiellement s'échapper d'un  optimum local très superficiel en acceptant un  voisin de même énergie (ou très légèrement pire) q'il est dans le top 5,mais finit invariablement bloqué face à un optimum plus profond.
    \item \textbf{random\_walk :} Ne se bloque jamais. Ignorant la topologie de la fonction coù, l'algorithme ressort naturellement de n'importe quel optimum local en marchant au hasard.
\end{itemize}
\end{answers}


\end{document}
